\begin{abstract}
Research on screen-based interactive systems takes months to stand up. Recruitment, pilot sessions, IRB, and the first three rounds of peer review each expose design flaws the authors could have caught earlier. We describe \textbf{Probe}, a command-line tool that runs a rough HCI research premise through an eight-stage pipeline of Claude Opus 4.7 and Sonnet 4.6 agents: premise interrogation, divergent ideation, literature grounding, Wizard-of-Oz specification, simulated walkthrough, capture-risk audit, adversarial review, and guidebook assembly. Each branch runs in a real git worktree; each stage's output validates against a JSON schema with one repair pass; every paragraph of the final guidebook carries a provenance tag that a markdown-AST linter enforces. A Managed Agents extension lets a reviewer agent run bash and grep against the branch artifacts to produce measurements a single-prompt audit cannot. On one demo premise, the adversarial methodologist caught an announcement-duration confound the authors missed; on an ablation trial, Sonnet 4.6 caught the same confound for 2.4\texttimes{} less cost but recommended reject where Opus recommended major revision. A planted fake citation did not produce a fabricated \texttt{[SOURCE\_CARD]} tag. The tool refuses to cross the rehearsal-to-evidence boundary through a combination of required provenance tags, a forbidden-phrase linter, and a dedicated \texttt{[HUMAN\_REQUIRED]} next-step block. Probe is not a replacement for user research. It is a triage layer between a rough premise and a study that a researcher can actually defend.
\end{abstract}
